\section{Discrete Formulation of Alignment Optimization}

\subsection{Discretization of the Domain}

\begin{definition}[Discretization]
Let
\[
0 = s_0 < s_1 < \dots < s_N = L
\]
be a partition of the interval \([0,L]\).
\end{definition}

\begin{remark}
We denote
\[
\Delta s_i = s_{i+1} - s_i.
\]
\end{remark}

\subsection{Discrete Variables}

\begin{definition}[Discrete curvature]
Let
\[
\kappa_i \approx \kappa(s_i)
\]
denote the curvature values at the nodes.
\end{definition}

\begin{definition}[Discrete cant]
Let
\[
\phi_i \approx \phi(s_i)
\]
denote the cant values.
\end{definition}

\begin{definition}[Parameter vector]
The optimization variable is
\[
\mathbf{x}
=
(\kappa_0,\dots,\kappa_N,\phi_0,\dots,\phi_N) \in \R^{2(N+1)}.
\]
\end{definition}

\subsection{Discrete Pose Reconstruction}

\begin{remark}
The pose3 state is reconstructed by numerical integration.
\end{remark}

\begin{definition}[Discrete integration]
Given initial conditions
\[
\gamma_0,\ \theta_0,\ \phi_0,
\]
we compute iteratively:
\[
\theta_{i+1} = \theta_i + \kappa_i \Delta s_i,
\]
\[
\gamma_{i+1} = \gamma_i + \Delta s_i
(\cos\theta_i,\sin\theta_i),
\]
\[
\phi_{i+1} = \phi_i + \omega_i \Delta s_i.
\]
\end{definition}

\begin{remark}
Here, \(\omega_i\) may be expressed via differences of \(\phi_i\).
This discrete reconstruction corresponds to standard numerical
integration ideas \cite{Hairer1993ODE,Press2007NumericalRecipes}.
\end{remark}

\subsection{Finite Differences}

\begin{definition}[First derivative]
Approximate derivatives by
\[
\kappa'(s_i)
\approx
\frac{\kappa_{i+1}-\kappa_i}{\Delta s_i},
\quad
\phi'(s_i)
\approx
\frac{\phi_{i+1}-\phi_i}{\Delta s_i}.
\]
\end{definition}

\begin{definition}[Second derivative]
\[
\kappa''(s_i)
\approx
\frac{\kappa_{i+1}-2\kappa_i+\kappa_{i-1}}{(\Delta s_i)^2}.
\]
\end{definition}

\subsection{Discrete Objective Functions}

\subsubsection{Curvature Smoothness}

\[
J_\kappa(\mathbf{x})
=
\sum_{i=0}^{N-1}
\left(
\frac{\kappa_{i+1}-\kappa_i}{\Delta s_i}
\right)^2
\Delta s_i.
\]

\subsubsection{Cant Smoothness}

\[
J_\phi(\mathbf{x})
=
\sum_{i=0}^{N-1}
\left(
\frac{\phi_{i+1}-\phi_i}{\Delta s_i}
\right)^2
\Delta s_i.
\]

\subsubsection{Dynamic Objective}

\[
J_{\mathrm{dyn}}(\mathbf{x})
=
\sum_{i=0}^{N}
\bigl(
v^2 \kappa_i - g \sin\phi_i
\bigr)^2
\Delta s_i.
\]

\subsubsection{Schwerpunkt Objective}

\begin{remark}
Using finite differences, we approximate
\[
\|\gamma_c''(s_i)\|^2
\]
by second-order differences of \(\gamma_c\).
\end{remark}

\[
J_{\mathrm{veh}}(\mathbf{x})
=
\sum_{i=1}^{N-1}
\|\gamma_{c,i+1} - 2\gamma_{c,i} + \gamma_{c,i-1}\|^2.
\]

\subsection{Combined Discrete Problem}

\begin{definition}[Discrete objective]
\[
J(\mathbf{x})
=
\alpha J_{\mathrm{veh}}
+
\beta J_{\mathrm{dyn}}
+
\gamma J_\kappa
+
\delta J_\phi.
\]
\end{definition}

\subsection{Constraints}

\begin{definition}[Equality constraints]
\[
\kappa_0 = 0,
\quad
\kappa_N = \kappa_1.
\]
\end{definition}

\begin{definition}[Inequality constraints]
\[
|\kappa_i| \le \kappa_{\max},
\quad
|\phi_i| \le \phi_{\max}.
\]
\end{definition}

\begin{remark}
Further constraints may include:
\begin{itemize}
\item bounds on derivatives,
\item continuity conditions,
\item alignment anchoring points.
\end{itemize}
\end{remark}

\subsection{Final Optimization Problem}

\begin{definition}[Discrete optimization problem]
Find
\[
\mathbf{x}^\ast \in \R^{2(N+1)}
\]
such that
\[
J(\mathbf{x}^\ast)
=
\min_{\mathbf{x}} J(\mathbf{x})
\]
subject to the given constraints.
\end{definition}

\subsection{Structure of the Problem}

\begin{remark}
The problem has the structure of a nonlinear least-squares problem.
\end{remark}

\begin{remark}
It is therefore well suited for:
\begin{itemize}
\item Sequential Quadratic Programming,
\item Gauss--Newton methods,
\item sparse optimization techniques.
\end{itemize}
\end{remark}

\subsection{Connection to Implementation}

\begin{remark}
The discrete formulation directly corresponds to a computational model:
\begin{itemize}
\item \(\kappa_i\) correspond to element parameters,
\item \(\phi_i\) correspond to cant samples,
\item reconstruction matches numerical evaluation in software.
\end{itemize}
\end{remark}

\begin{summary}
The continuous alignment optimization problem can be reduced to a finite
dimensional nonlinear optimization problem.

This provides the direct link between theoretical modelling and
practical implementation.
\end{summary}

